% Generated draft: requires mathematical and source review.
\documentclass[11pt]{article}
\usepackage[margin=25mm]{geometry}
\usepackage{fontspec}
\setmainfont{Latin Modern Roman}
\usepackage{amsmath,amssymb,mathtools,unicode-math}
\setmathfont{Latin Modern Math}
\usepackage{enumitem,needspace,xurl,hyperref}
\setlength{\parindent}{0pt}
\setlength{\parskip}{5pt}
\setlength{\emergencystretch}{3em}
\newcommand{\N}{\mathbb{N}}
\newcommand{\Z}{\mathbb{Z}}
\newcommand{\Q}{\mathbb{Q}}
\newcommand{\R}{\mathbb{R}}
\newcommand{\C}{\mathbb{C}}
\newcommand{\F}{\mathbb{F}}
\newcommand{\A}{\mathbb{A}}
\newcommand{\PP}{\mathbb{P}}
\newcommand{\E}{\mathbb{E}}
\newcommand{\eps}{\varepsilon}
\newcommand{\dd}{\,\mathrm{d}}
\begin{document}
\section*{UnsolvedMath ID: 700375}
\textbf{Source problem code:} catalogue entry
\subsection{Definitions and mathematical statement}
For a topological group \(G\), let
\[
\operatorname{Aut}_{\mathrm{top}}(G)=\{\alpha:G\to G:\alpha\text{ is a group isomorphism and both }\alpha,\alpha^{-1}\text{ are continuous}\},
\]
with composition as the group operation. For an abstract group \(H\), let \(\operatorname{Aut}(H)\) denote its group of all group automorphisms. When \(H\) carries the discrete topology, every automorphism is a topological automorphism.

The question proposed by O. V. Mel'nikov is whether, for every finite group \(A\),
\[
\left[\exists\text{ a locally compact topological group }G\quad \operatorname{Aut}_{\mathrm{top}}(G)\cong A\right]
\Longrightarrow
\left[\exists\text{ a group }H\quad \operatorname{Aut}(H)\cong A\right].
\]
All displayed isomorphisms are isomorphisms of abstract groups; the record specifies no topology on an automorphism group. The group \(H\) is to be given the discrete topology, and no restriction on its cardinality is stated.

The phrase ``locally compact group'' retains the convention of the cited source: the supplied record does not explicitly specify whether Hausdorffness is included, so that separation convention requires confirmation before the scope is fixed.

The supplied statement reports that the implication holds when \(A\) is cyclic and that local compactness of \(G\) is essential, citing R. J. Wille. These are supplied background assertions, not independently verified conclusions of this draft; no counterexample supporting the latter assertion is included in the record.
\subsection{Short English statement}
If a finite group \(A\) occurs as the full group of topological automorphisms of some locally compact group, must it also occur as the full automorphism group of some discrete group? The discrete group need not be finite. The source reports an affirmative answer for cyclic \(A\) and says that local compactness is essential; its separation convention needs confirmation.
\subsection{Sources}
\begin{itemize}
\item Local UnsolvedMath registry, permanent ID \(700375\), exact source problem code \texttt{catalogue entry}; supplied complete source record, upstream revision \texttt{d371ee917ac9647fbb06cd6fc529deec6f3dda82}. The registry labels this problem \texttt{open}; that label is supplied metadata, not an independently verified literature assessment.
\item O. V. Mel'nikov, Problem 6.29, \emph{The Kourovka Notebook (Unsolved Problems in Group Theory)}, originating in the 6th issue (1978), as identified in the supplied record with the 21st issue edition (2026), arXiv:\texttt{1401.0300v48}. The supplied collection description names E. I. Khukhro and V. D. Mazurov as editors.
\item R. J. Wille, \emph{Indag. Math.}, \textbf{25}, no. 2 (1963), 218--224. This reference is supplied parenthetically after the assertions about cyclic \(A\) and the necessity of local compactness. The record supplies no article title or further bibliographic details.
\end{itemize}
\end{document}
